% Drawing 78/05-5 -- folder 78, batch 05, archive page 100 (PDF page 101).
% Transcription: à gauche, un croquis : les axes $x$, $y$ et les axes obliques $v$, $w$ par l'origine.
% Kind: axes x, y and oblique axes v, w. Method 2 (schematic redraw from the facsimile). Status: ok.
% Notes: The axes x, y and the oblique axes v, w through the origin; v drawn heavier.
% Coordinates: 1 unit = 1 mm = 1/100 of the crop width (bbox in page fractions: [0.06, 0.3, 0.32, 0.46]).
% Redrawn by Claude (claude-opus-5-5), issue #73, 2026-11-07; not checked by a human; not compiled (no TeX in the session).
% Rights: reproduces part of the work -- behind RENDER_TRANSCRIPTS, never published by itself.
\documentclass[tikz,border=3pt]{standalone}
\usepackage{amsmath,amssymb}
\usetikzlibrary{patterns,arrows.meta}
\tikzset{pencil/.style={black!45}, struck/.style={black!35}, mark/.style={circle,draw,inner sep=1.3pt}, pt/.style={circle,fill,inner sep=0.9pt}, opt/.style={circle,draw,fill=white,inner sep=1.1pt}}
\begin{document}
\begin{tikzpicture}[x=1mm,y=1mm,line cap=round,line join=round,font=\scriptsize]
\draw (18,-52) -- (86,-50);
\draw (42,-24) -- (42,-75);
\draw (26,-67) -- (67,-21);
\draw (25,-38) -- (41,-51);
\draw[very thick] (41,-51) -- (76,-77);
\node[inner sep=1pt] at (94,-50) {\scriptsize $x$};
\node[inner sep=1pt] at (40,-17) {\scriptsize $y$};
\node[inner sep=1pt] at (70,-18) {\scriptsize $w$};
\node[inner sep=1pt] at (82,-80) {\scriptsize $v$};
\end{tikzpicture}
\end{document}
